%==============================================================================
% results_kscaling_eicu.tex — external validation of the records-per-individual
% law on eICU-CRD. MERGED: \input from results.tex after "Records per
% individual" (sec:results:kscaling).
%
% Every number below is read from an artifact:
%   results/phase5/k_scaling/eicu/k{K}_seed{S}/hist_{arm}_test.json
%       -> auprc (primary), macro_auc (secondary)
%   results/phase5/k_scaling/eicu/k_scaling_eicu.csv   (collected summary)
%   document/01_planning/PREREGISTRATION_EICU.md       (frozen forecast, §2;
%                                                       amendment A1)
%   data/processed/eicu/meta/dataset_meta.json         (corpus description)
%
% k = 16 COMPLETE at three seeds as of 21 September 2026; the row is in the
% table and P1c is settled below. Per-seed gaps +0.0146, +0.0222, -0.0054.
%==============================================================================

\subsubsection{External validation of the records-per-individual law}
\label{sec:results:kscaling:eicu}

The law of Section~\ref{sec:results:kscaling} was measured on a single corpus
whose clients are surgical specialties rather than institutions, so it is not
established whether it characterizes patient-level private federated learning
in general or only that cohort. It is therefore tested on a second clinical
corpus of different construction: $3{,}782{,}805$ hourly windows from
$132{,}736$ individuals across $187$ clients that are \emph{real hospitals},
with a median of $444$ individuals per client and $28.50$ records per
individual. The label is in-unit mortality, with a prevalence of $0.0422$.

Two properties of this replication are relevant. First, its encoders were
pretrained on $526{,}867$ windows verified to be disjoint from every client's
training data, so the results do not carry the oracle-configuration caveat that
applies elsewhere in this paper. Second, the prediction was fixed before the
corpus was constructed. Fitting the four VitalDB points of
Table~\ref{tab:kscaling} gives
\begin{equation}
\mathrm{gap}(k) = A\left(1 - e^{-b(k-1)}\right),
\label{eq:kcurve}
\end{equation}
with $A = 0.0878 \pm 0.0035$ and $b = 0.883 \pm 0.129$ at a root-mean-square
error of $0.0028$. This curve, together with an acceptance band of $\pm 0.030$
derived from the VitalDB seed spread, was registered before any record of the
second corpus was accessed.

\begin{table}[!t]
\centering
\caption{Average precision by records per individual $k$ on the external
corpus, three seeds, paired by seed. Intervals are Student-$t$ half-widths as
elsewhere. The row $k = 16$ lies beyond the range of the first corpus. The final column is
the value predicted by Eq.~\eqref{eq:kcurve}, fixed before the corpus was
built.}
\label{tab:kscaling_eicu}
\setlength{\tabcolsep}{4pt}
\begin{tabular}{cccccc}
\toprule
$k$ & Record & Individual & Gap & $p$ & Predicted \\
\midrule
$1$ & $0.2105$ & $0.2212$ & $-0.0107 \pm 0.0211$ & $0.1594$ & $0.0000$ \\
$2$ & $0.2245$ & $0.2185$ & $+0.0060 \pm 0.0100$ & $0.1215$ & $0.0515$ \\
$4$ & $0.2318$ & $0.2208$ & $+0.0110 \pm 0.0714$ & $0.5746$ & $0.0816$ \\
$8$ & $0.2374$ & $0.2239$ & $+0.0135 \pm 0.0592$ & $0.4296$ & $0.0876$ \\
$16$ & $0.2363$ & $0.2258$ & $+0.0105 \pm 0.0354$ & $0.3310$ & $0.0878$ \\
\bottomrule
\end{tabular}
\end{table}

The control holds. At $k = 1$, the gap is $-0.0107$, with an interval covering
zero and a magnitude within the pre-specified tolerance of $0.020$. The two arms
are therefore not separated where the construction requires them to agree, and
the remaining rows can be attributed to the privacy unit.

The point forecasts are not confirmed. At every $k > 1$, the gap lies outside
the $\pm 0.030$ band, and the measured cost is approximately one sixth of the
predicted cost: $0.0135$ against $0.0876$ at $k = 8$.

Two properties of the law are reproduced. The first is the ordering: the gap is
negative at $k = 1$, positive at $k = 2$, and increases through $k = 4$ to
$k = 8$, as described by Eq.~\eqref{eq:kcurve}. This is consistent with the
mechanism, since the bound on an individual's summed contribution is tightest
for individuals who contribute additional records.

The second is the saturation, which the first corpus cannot test. The
registered prediction was that the gap beyond $k = 8$ would exceed the gap at
$k = 8$ by less than $0.015$. Across paired seeds, the increment from $k = 8$ to
$k = 16$ is $-0.0030$, with an interval of $[-0.0316, +0.0255]$ ($p = 0.692$).
The gap does not continue to increase, and the registered bound is met. The
cost of the guarantee therefore approaches a ceiling in $k$ on both corpora.

The asymmetry between the arms is also reproduced. Across all $k$, the
individual-level arm varies within $[0.2185, 0.2258]$, a range of $0.0073$
against the pre-specified $0.020$, whereas the record-level arm increases from
$0.2105$ to $0.2374$. As on the first corpus, the widening gap reflects
improvement of the record-level arm rather than degradation of the
individual-level arm. The corpora therefore differ in the magnitude of the
effect but not in its mechanism, and, as registered in advance, the shape is
interpreted as structural and the magnitude as corpus-dependent.

The statistical power of this replication is limited. At three seeds, no gap in
Table~\ref{tab:kscaling_eicu} is resolved: every interval covers zero, and no
$p$-value is below $0.12$. The seed standard deviation of the gap at $k = 8$ is
of the same order as the gap, and resolving an effect of this size at
conventional power would require approximately fifty seeds, against the three
fixed in advance. The seed count was not increased after inspecting the
results. The trend is therefore reported as directionally consistent with the
VitalDB law and unresolved on this corpus. The area under the ROC curve is
less informative: its gaps over the same values of $k$ are $-0.0012$,
$-0.0015$, $+0.0018$, $+0.0001$ and $+0.0105$, with no trend over the range
covered by the first corpus. Average precision is the primary measure on this
corpus, as on the first, and the area under the ROC curve provides no evidence
either way.

One possible explanation for a preserved shape at a lower ceiling is offered as
a hypothesis. The two tasks differ in attainable precision by approximately a
factor of two: the arms on this corpus lie within $[0.1960, 0.2446]$ at a
prevalence of $0.0422$, whereas those on the first corpus lie between $0.4967$
and $0.5955$. A gap in average precision is bounded by the attainable
precision, so a cost that is a fixed fraction of attainable performance would
produce a smaller absolute gap on the harder task with the same ordering in
$k$. Under this explanation, the transferable quantity is the cost relative to
attainable precision rather than the ceiling $A$. Distinguishing this from a
corpus-specific ceiling requires a third corpus.

The range of $k$ is limited by the corpus. The registered design specified $k$
up to $32$. On this corpus, however, a budget of twenty-four windows per
admission means that only $20.2\,\%$ of individuals reach thirty-two records,
and these individuals have multiple admissions and are not representative of
the cohort. Estimating the arm at $k = 32$ would therefore measure the effect
of readmission rather than of the privacy unit. The maximum was reduced to
$16$, which retains $91.0\,\%$ of individuals, and $k = 32$ is reported as
untested. This change was recorded as a dated amendment to the registration
before any result beyond $k = 8$ was computed, and it did not alter any
predicted value.

%==============================================================================
% COMPLETE. k = 16 landed 21 September 2026 at three seeds; P1c is settled in
% the body above. No pending artifact remains for this subsection.
%   artifacts : results/phase5/k_scaling/eicu/k_scaling_eicu.csv (30 rows)
%               results/phase5/k_scaling/eicu/k{1,2,4,8,16}_seed{42,43,44}/
%   k = 32 is reported UNTESTED, not dropped — amendment A1.
%==============================================================================
